\section{Konsep Interupsi dan Tabel Vektor}
\label{sec:konsep-interupsi}

Pada Bab~\ref{chap:pengenalan}, layanan \code{INT 21h} dipakai sebagai
kotak hitam. Bab ini membuka kotak hitam tersebut: pertanyaannya bukan
lagi ``bagaimana memakai'' layanan, melainkan ``apa yang sesungguhnya
terjadi'' ketika \asm{INT} dijalankan.

\subsection{Mengapa Interupsi Diperlukan}

Prosesor berjalan jauh lebih cepat daripada perangkat di sekitarnya.
Bila program terus-menerus memeriksa status keyboard, sebagian besar
waktunya terbuang untuk pengulangan yang hampir selalu sia-sia. Cara
ini disebut \eng{polling}.

\begin{definisi}
\term{Interupsi} adalah mekanisme yang memaksa prosesor menghentikan
sementara alur program, memindahkan kendali ke rutin khusus yang
disebut \term{penangan interupsi} (\eng{interrupt handler}), lalu
melanjutkan kembali program semula tepat pada instruksi berikutnya.
\end{definisi}

Interupsi juga menjadi sarana baku untuk meminta layanan sistem:
program cukup memanggil \code{INT 10h} tanpa menulis sendiri rutin
pengendali video \cite{intel1979}.

\subsection{Tiga Jenis Interupsi}

Berdasarkan sumber pemicunya, interupsi dibagi menjadi tiga jenis.

\begin{enumerate}
    \item \term{Interupsi perangkat keras} berasal dari luar prosesor
          melalui kaki \file{INTR} atau \file{NMI}. Interupsi melalui
          \file{INTR} bersifat \eng{maskable}: dapat ditahan bila
          \flag{IF} bernilai nol. Interupsi melalui \file{NMI} tidak
          dapat ditolak dan dipakai untuk kejadian gawat.

    \item \term{Interupsi perangkat lunak} dipicu instruksi
          \lstasm|INT n| di dalam program. Instruksi ini sengaja
          dijalankan untuk meminta layanan; seluruh layanan BIOS dan
          DOS termasuk jenis ini.

    \item \term{Pengecualian} (\eng{exception}) dibangkitkan prosesor
          ketika terjadi keadaan yang tidak diinginkan. Instruksi
          \lstasm|DIV| dengan pembagi nol memicu vektor \hex{00},
          sedangkan \lstasm|INTO| memicu vektor \hex{04} bila
          \flag{OF} bernilai satu.
\end{enumerate}

Perbedaan ketiga jenis tersebut dirangkum pada
Tabel~\ref{tab:jenis-interupsi}.

\begin{table}[H]
    \centering
    \caption{Tiga jenis interupsi pada prosesor 8086.}
    \label{tab:jenis-interupsi}
    \small
    \begin{tabular}{@{}p{2.5cm}p{3.1cm}p{4.1cm}p{3.9cm}@{}}
        \toprule
        \textbf{Jenis} & \textbf{Sumber} & \textbf{Contoh} &
        \textbf{Sifat} \\
        \midrule
        Perangkat keras & Kaki \file{INTR} atau \file{NMI} &
        Timer 18,2 kali per detik (\hex{08}); penekanan tombol
        (\hex{09}) & Terjadi kapan saja, tidak direncanakan program \\
        \addlinespace
        Perangkat lunak & Instruksi \lstasm|INT n| &
        \code{INT 21h} fungsi \hex{09} untuk mencetak string &
        Disengaja, hasilnya dapat diperkirakan \\
        \addlinespace
        Pengecualian & Prosesor, akibat kesalahan eksekusi &
        Pembagian dengan nol (\hex{00}); limpahan \lstasm|INTO|
        (\hex{04}) & Menandakan kesalahan yang harus ditangani \\
        \bottomrule
    \end{tabular}
\end{table}

\subsection{Tabel Vektor Interupsi}

Prosesor harus mengetahui letak rutin penangan setiap interupsi. Untuk
itu, 256 kemungkinan interupsi pada 8086 diberi nomor \hex{00} sampai
\hex{FF}, dan alamat setiap penangan disimpan pada \term{tabel vektor
interupsi} (\eng{interrupt vector table}). Tabel ini menempati 1.024
byte pertama memori, yaitu alamat \hex{0000} sampai \hex{03FF} pada
segmen nol, dan mampu menampung $1.024 \div 4 = 256$ vektor.

Setiap entri terdiri atas 2 byte \term{offset} penangan yang dimuat ke
\reg{IP}, diikuti 2 byte \term{segmen} yang dimuat ke \reg{CS}. Bagian
bernilai lebih rendah diletakkan pada alamat yang lebih rendah. Tata
letaknya tampak pada Gambar~\ref{fig:tabel-vektor}.

\begin{figure}[H]
    \centering
    \begin{tikzpicture}[
        alamat/.style={font=\ttfamily\scriptsize, text=obeGray,
                       anchor=east},
        sel/.style={draw=obeGray, line width=0.6pt,
                    minimum height=0.5cm, minimum width=1.45cm,
                    inner sep=0pt, anchor=west},
        selisi/.style={sel, fill=obeBlueLight},
        ket/.style={font=\scriptsize, anchor=west, text=obeBlue},
        kopf/.style={font=\scriptsize\bfseries, text=obeGray}
    ]
        \node[kopf, anchor=east] at (1.55,0) {vektor};
        \node[kopf, anchor=east] at (3.05,0) {alamat};
        \node[kopf, anchor=west] at (3.30,0) {isi 4 byte};
        \node[kopf, anchor=west] at (6.45,0) {keterangan};
        \draw[obeGray, line width=0.5pt] (-1.7,0.28) -- (9.4,0.28);

        \foreach \r/\v/\a/\k in {%
            0/00h/0000h/{pembagian dengan nol},
            1/01h/0004h/{eksekusi satu langkah},
            2/02h/0008h/{interupsi NMI},
            3/08h/0020h/{timer sistem},
            4/09h/0024h/{keyboard},
            5/10h/0040h/{layanan video BIOS},
            6/16h/0058h/{layanan keyboard BIOS},
            7/17h/005Ch/{layanan printer BIOS},
            8/1Ch/0070h/{timer tick bagi pengguna},
            9/21h/0084h/{layanan DOS},
            10/FFh/03FCh/{belum ditetapkan}%
        }{%
            \pgfmathsetmacro{\y}{-0.58*\r-0.30}
            \node[alamat] at (1.55,\y) {\v};
            \node[alamat] at (3.05,\y) {\a};
            \node[selisi] at (3.30,\y) {offset};
            \node[selisi] at (4.75,\y) {segmen};
            \node[ket] at (6.45,\y) {\k};
        }
    \end{tikzpicture}
    \caption{Tata letak tabel vektor interupsi pada segmen nol, dari
    alamat \hex{0000} sampai \hex{03FF}. Entri berisi alamat penangan,
    bukan penangan itu sendiri.}
    \label{fig:tabel-vektor}
\end{figure}

\subsection{Menghitung Alamat Vektor ke-\texorpdfstring{$n$}{n}}

Karena setiap vektor berukuran tetap, alamat awal vektor bernomor $n$
dihitung dengan rumus berikut:

\begin{equation}
\label{eq:alamat-vektor}
\text{alamat vektor}_n = n \times 4
\end{equation}

Hasilnya adalah offset di dalam segmen nol. Vektor \hex{09} (keyboard)
berada pada $9 \times 4 = 36 = \hex{24}$, vektor \hex{10} (video BIOS)
pada $16 \times 4 = 64 = \hex{40}$, dan vektor \hex{21} (layanan DOS)
pada $33 \times 4 = 132 = \hex{84}$. Vektor yang sering dijumpai pada
praktikum dirangkum pada Tabel~\ref{tab:vektor-penting}.

\begin{table}[H]
    \centering
    \caption{Vektor interupsi penting beserta alamat dan kegunaannya.}
    \label{tab:vektor-penting}
    \small
    \begin{tabular}{@{}ccp{3.0cm}p{6.0cm}@{}}
        \toprule
        \textbf{Vektor} & \textbf{Alamat entri} & \textbf{Pemicu} &
        \textbf{Kegunaan} \\
        \midrule
        \hex{00} & \hex{0000} & Pengecualian prosesor &
        Pembagian dengan nol, dipicu otomatis oleh \lstasm|DIV| \\
        \addlinespace
        \hex{01} & \hex{0004} & Pengecualian prosesor &
        Eksekusi satu langkah, dipicu saat \flag{TF} bernilai satu \\
        \addlinespace
        \hex{08} & \hex{0020} & Perangkat keras &
        Timer sistem, dibangkitkan 18,2 kali per detik \\
        \addlinespace
        \hex{09} & \hex{0024} & Perangkat keras &
        Penekanan dan pelepasan tombol keyboard \\
        \addlinespace
        \hex{10} & \hex{0040} & Perangkat lunak &
        Layanan video BIOS: mode layar, kursor, penulisan karakter \\
        \addlinespace
        \hex{16} & \hex{0058} & Perangkat lunak &
        Layanan keyboard BIOS: membaca dan memeriksa tombol \\
        \addlinespace
        \hex{17} & \hex{005C} & Perangkat lunak &
        Layanan printer BIOS \\
        \addlinespace
        \hex{1A} & \hex{0068} & Perangkat lunak &
        Waktu dan tanggal dari jam perangkat keras \\
        \addlinespace
        \hex{1C} & \hex{0070} & Perangkat lunak &
        Timer \eng{tick} bagi pengguna, dipanggil dari penangan
        \hex{08} \\
        \addlinespace
        \hex{21} & \hex{0084} & Perangkat lunak &
        Layanan DOS: masukan, keluaran, berkas, terminasi program \\
        \bottomrule
    \end{tabular}
\end{table}

\subsection{Mekanisme \asm{INT} dan \asm{IRET}}

Instruksi \lstasm|INT n| bekerja seperti pemanggilan prosedur jarak
jauh, tetapi dengan tambahan penyimpanan \eng{flag}
\cite{intel1979}. Urutannya:

\begin{enumerate}
    \item isi register \reg{FLAGS} disimpan ke stack;
    \item register \reg{CS} disimpan ke stack;
    \item register \reg{IP} --- yang ketika itu sudah menunjuk
          instruksi sesudah \lstasm|INT n| --- disimpan ke stack;
    \item \flag{IF} dan \flag{TF} dikosongkan, sehingga interupsi
          perangkat keras ditahan sementara dan eksekusi satu langkah
          tidak aktif;
    \item alamat penangan dibaca dari entri vektor pada alamat
          $n \times 4$; dua byte pertama dimuat ke \reg{IP} dan dua
          byte berikutnya ke \reg{CS}.
\end{enumerate}

Ketiga penyimpanan itu menurunkan \reg{SP} sebanyak 6 byte. Isi stack
sesudahnya tampak pada Gambar~\ref{fig:stack-int}.

\begin{figure}[H]
    \centering
    \begin{tikzpicture}[
        sel/.style={draw=obeTeal, line width=0.7pt, fill=obeTealLight,
                    minimum width=3.0cm, minimum height=0.55cm,
                    font=\ttfamily\small, inner sep=0pt},
        old/.style={draw=obeGray, line width=0.7pt, fill=obeGrayLight,
                    minimum width=3.0cm, minimum height=0.55cm,
                    font=\ttfamily\small, inner sep=0pt},
        alamat/.style={font=\ttfamily\scriptsize, text=obeGray,
                       anchor=west},
        panah/.style={-{Stealth[length=2.2mm]}, line width=0.9pt,
                      obeMaroon}
    ]
        \node[sel] (ip) at (0,0) {\reg{IP}};
        \node[sel] (cs) [below=0pt of ip] {\reg{CS}};
        \node[sel] (fl) [below=0pt of cs] {\reg{FLAGS}};
        \node[old] (lama) [below=0pt of fl] {data lama};

        \node[alamat] at (1.65,0) {SS:SP};
        \node[alamat] at (1.65,-0.55) {SS:SP+2};
        \node[alamat] at (1.65,-1.10) {SS:SP+4};

        \draw[panah] (-2.9,0) -- (-1.62,0);
        \node[font=\scriptsize, text=obeMaroon, anchor=east]
            at (-3.0,0) {SP sesudah \asm{INT}};

        \node[font=\scriptsize, text=obeGray, anchor=north west,
              align=left] at (-3.0,-2.35)
            {SP berkurang 6 byte. Pengambilan\\kembali oleh \asm{IRET}
             dilakukan\\dalam urutan terbalik.};
    \end{tikzpicture}
    \caption{Isi stack sesudah instruksi \asm{INT n} dijalankan,
    dengan \reg{SP} menunjuk pada nilai \reg{IP} yang tersimpan paling
    akhir.}
    \label{fig:stack-int}
\end{figure}

Instruksi \lstasm|IRET| melakukan tiga pengambilan berurutan, yaitu
\reg{IP}, \reg{CS}, lalu \reg{FLAGS}, persis kebalikan urutan
penyimpanan. Bandingkan dengan \lstasm|CALL| pada
Bab~\ref{chap:prosedur} yang hanya menyimpan \reg{CS} dan \reg{IP}:
\lstasm|IRET| setara dengan \lstasm|RET| ditambah \lstasm|POPF|.

\begin{peringatanbox}
Rutin penangan interupsi wajib diakhiri \lstasm|IRET|, bukan
\lstasm|RET|. Bila \lstasm|RET| dipakai, dua byte teratas stack yang
diambil sebagai alamat kembali justru berisi \reg{FLAGS}, sehingga
prosesor melompat ke alamat acak dan program menghentikan diri atau
sistem membeku.
\end{peringatanbox}

\subsection{Flag \flag{TF} dan Eksekusi Satu Langkah}

Bit ke-8 register \reg{FLAGS} adalah \term{trap flag} \flag{TF}. Bila
bit ini bernilai satu, prosesor menjalankan satu instruksi lalu segera
membangkitkan interupsi vektor \hex{01}. Mekanisme inilah yang dipakai
\eng{debugger} untuk menyediakan tombol \file{Single Step} pada Emu8086
atau perintah \file{T} pada DEBUG \cite{ebrahim2018}.

Peran \flag{TF} menjelaskan langkah keempat mekanisme \lstasm|INT| di
atas: bila flag itu tetap bernilai satu, penangan akan memicu vektor
\hex{01} setiap kali satu instruksinya selesai sehingga prosesor
terjebak. Pemakaiannya pada penelusuran program dibahas pada
Bab~\ref{chap:debugging}.
